\chapter{Den tidlige historie}
Dette afsnit er en gengivelse af det tilsvarende afsnit i J50 skriftet. Da resten af J70 skriftet har fokus på årene 2016 til 2026, har vi valgt at lade denne oprindelige fortælling om \TKETs historie forekomme næstent uden ændringer: det eneste der er tilføjet er afsnitsinddelinger for at gøre det mere læsevenligt i den trykte version af dette skrift.

\section{\TKETs historie}
\vspace{-0.2cm}
\textit{\textbf{---} en abrupt men saglig gennemgang}
\newline\newline
\subsection{1950'erne: \TKETs start}
\noindent Den 11. oktober 1956 var der en gruppe fremsynede unge mennesker fra det naturvidenskabelige fakultet ved Aarhus Universitet, der besluttede, at stedet manglede en foredrags- og festforening. De dannede derfor en sådan, der fik en bestyrelse på fem personer: Formand, næstformand, kasserer, ceremonimester og sekretær. Denne forening blev oprindeligt stiftet som søsterforening til „Parentesen“ ved Københavns Universitets Naturvidenskabelige Fakultet. Dog var man lokalpatrioter idet der blev lavet en aftale med Ceres om levering af øl 3–4 gange årligt! Altså ingen Hof eller Tuborg i Århus. Vi er ikke længere helt så loyale, men der arrangeres da årligt en udflugt ned til bryggeriet, hvor vi får en rundvisning. Det første stykke tid havde foreningen ikke et reelt navn, men på generalforsamlingen i 1957 krævede sekretæren hårdnakket at et blev fundet, og efter lang tids diskussion blev man enige om navnet \TKET. Således blev det fastlagt hvad den lille forening skulle hedde, og den var nu i sandhed på vej mod ære og berømmelse.

I starten var det kun T’et der altid blev skrevet med stort, men i 1958 ser vi første gang at hele navnet skrives med stort. Det er fra året 1976–1977 vi har det første skriftlige bevis for at \TKET ikke må deles. Derimod er traditionen med at \TKET skal skrives med stort opstået i tidsrummet maj 1985 til oktober 1985, hvor det eksplicit står som regel i en officiel skrivelse. Det præcise tidspunkt har ikke kunnet findes.

\subsection{1960'erne: Foredrag, fester og arrangementer}
De første år drejede det sig mest om at holde foredrag, men man besluttede dog hurtigt at afholde en stiftelsesfest én gang årligt, hvor ceremonimesteren officielt skulle indsættes i sit hverv. Det var således i 1963 at ceremonimesterkæden for første gang blev fremvist for en måbende folkemængde. Men i november 1970 bliver kæden smidt væk i en brandert, og i stedet må \CERM indvies med en belgisk studenterhue. Året efter var den nye og nuværende kæde klar. Julefesten går også tilbage til de tidligste annaler, men i de tidlige år indledte man med et foredrag, og det var først i 1962 det blev tradition at indlede julefesten med en julerevy. Faste majrevyer kom til efterhånden. Det skal dog nævnes, at den første \TK-revy i historien blev opført i maj 1958.

De fester der efterhånden kom til i årenes løb var i øvrigt vidt berømte for den punch, der blev serveret. Der var tale om 1–2 mælkejunger (sådan nogle er \textit{store}) med punch, og den alkohol, der var i, var ikke for børn. Kemikerne var endnu ikke brudt ud for at starte \alkymia, så de var med i \TKET, og via dem skaffedes finsprit.

Men lad os nu vende tilbage til bestyrelsesåret 1963–1964. Ceremonimesterkæden var ikke det eneste bestyrelsen havde i ærmet. Den 6. februar besluttede de sig for at hyre en til at lave plakater og lignende til fester og foredrag, hvilket altså er første gang man har noget, der ligner en PR-titel i \TKET. I bestyrelsen 1964–1965 optræder der ingen PR-mand, men ved generalforsamlingen i 1965 valgtes der for første gang en syvpersonersbestyrelse. Dog er PR og \VC ikke fast titelgivet endnu, det sker først fra generalforsamlingen i 1970. Betegnelsen BEST går helt tilbage til 1958.

Som nævnt var \TKET i starten hovedsagligt en foredragsforening, og man afholdt da også en masse foredrag med talrige personer, der i dag er kendt for deres virke på fakultetet. Desuden kom der tit folk udefra; hovedsagligt fra Københavns Universitet, men også fra andre fakulteter, DTU, erhvervslivet og lignende.

I 1966 startede man en tradition, der puttede lidt kultur ind i de studerendes hoveder. Bestyrelsen havde nemlig købt to forestillinger på Svalegangen, og billetter til disse solgte man så til medlemmerne. Det hændte selvfølgelig ofte, at teatergængerne bagefter fandt et spise- eller værtshus, hvor de fortsatte med at hygge sig. Denne ellers udmærkede tradition er desværre med tiden forsvundet.

I sommeren 1967 skete der en skelsættende begivenhed for både fakultetet og \TKET. Man flyttede nemlig til de nye fakultetsbygninger på Ny Munkegade. \TKET blev straks en del af byggeplanerne med lokalet B1.28, hvor det som bekendt har ligget lige siden. I 2004–2005 var man vi nødt til midlertidigt at forlade vort elskede kammer pga. renovering.

\TKET har en stærk tradition for sport. Med begyndelse allerede i slutningen af 1960’erne har \TKET jævnligt arrangeret forskellige former for konkurrencer og spil, heriblandt whist, skak, bingo, fodbold og bordtennis. I samarbejde med Provokation (et blad der var startet under \TKET, og i en årrække trivedes på fakultetet ved at leve op til sit navn) arrangerede man vektorrally, også kaldet provorally og meget andet. I dag holdes de ædle sportsgrene i hævd ved den årlige TK-Open, der finder sted i Vandrehallen. De hyppigste discipliner er nu vektorrally, makkercasino, euchre (udtales: jugår) og det legendariske ølluffo.

\subsection{1970'erne: Spil og HSTR}
Bestyrelsen 1971–1972 skal især huskes for den store indsats de gjorde for de sportslige aktiviteter, idet de startede to arrangementer af sportslig karakter, der stadig lever i bedste velgående: Den første luttho-turnering med kæmpe finalebræt, øl-brikker og alt det der hører med til en rigtig lucio-tunering, og den første løv- faldstur. Det var dog først flere år senere at turens destination blev hemmeligholdt indtil man ankom. Det skal nævnes at man allerede i 1965 havde forsøgt sig med en sådan tur, men opgav idet der ingen ledige kroer var. Det afsluttende traktement bestående af biksemad og spejlæg med en enkelt snaps har været fast tradition lige fra starten. At løvfaldsturen siden har skiftet navn til „Hemmelig Skovtur Til Rønde“ eller HSTR skyldes, som mange ved, Carsten Hansen, der var FORM i 1985–1986, idet han troskyldigt fortalte FORMjuntaen hvor turen gik hen. Ellers kunne de jo ikke skrive den sang, som de bildte ham ind, der var tradition for at FORMjuntaen skrev.

HSTR er ikke det eneste ud i naturen-arrangement \TKET afholder. I 1978 iværksatte man den første cykeltur. Disse cykelture har ikke været et fast indslag gennem alle årene, men de dukker op fra tid til anden. Først fra 2003 er det blevet tradition at der årligt cykles til Kalø, hvor man griller, hygger og overnatter inden man cykler hjem igen.

\subsection{Løbesedler}
Hvis der er nogen, der spekulerer på hvor gamle NF’s løbesedler er, er den ældste, bevarede løbeseddel fra september 1957, og den har, som moderne løbesedler, det format der hedder A6. Det vides at disse løbesedler, fra engang i 1959, stort set alle er skrevet på \TKETs officielle skrivemaskine, \RemToR. Så de russer og andre stakler (NF?), der tror at \RemToR er en hypermoderne, multiavanceret word-processor-model tager grundigt fejl. Det må dog indrømmes at den trods alle de år den har på bagen skriver forbløffende godt. Men nu vi alligevel er ved NF og \RemToR, kan det da lige nævnes at underskriften NF/\RemToR første gang optrådte på en løbeseddel i slutningen af september 1980.

\subsection{Generalforsamling}
Valget til \TKETs bestyrelse har alle dage været demokratisk. Men sidenhen er vi begyndt at vælge bestyrelsen på \textit{demokratisk} vis. I 1974 skete der dog noget bemærkelsesværdigt. På generalforsamlingen blev der først valgt en FORM og derpå en \CERM og endelig resten af bestyrelsen. Intet nyt i det, sådan står der jo i vores vedtægter at vi skal. Det usædvanlige bestod i, at af de flere lister, der var opstillet, vandt den den gamle bestyrelse; altså et kup! Det viste sig dog, utroligt men sandt, at denne bestyrelse, der jo havde siddet i ét år allerede, ikke kunne samarbejde, og det eneste de kunne blive enige om var at trække sig igen. Derefter blev den næste rigtige bestyrelse valgt, og de kunne samarbejde. Denne procedure, i forskellige afskygninger, er blevet anvendt med jævne mellemrum lige siden.

\subsection{Stiftelsesfest}
Den 8. november 1974 var det blevet tid til indsættelse af \CERM ved den traditionelle stiftelsesfest. Her var \CERM i mange år blevet indsat ved en bestemt ceremoni, hvor han fik overrakt kæden og hornet, men det år blev der desuden læst op af en stor bog, hvorudaf støvet væltede. Det var jo en meget gammel bog, ja helt fra vikingetiden endda.

\subsection{Mænd i kjoler}
Den klassiske cancan ses i protokollerne første gang i 1969, men allerede dengang var det en gammel tradition, idet billedteksten lyder: „Og til allersidst naturligvis can-can“. En anden stolt juletradition er afsyngelsen af „Kold X-mas“. Men det var faktisk sådan at man oprindeligt sang originalen „White Christmas“. I hvert fald indtil 1981 hvor den nye tekst for første gang blev afsunget til megen morskab for publikum, og i dag er det vel en lige så stor tradition som cancanen er det. I 1991 holdt man for første gang luciaoptog gennem institutternes gange, auditorier og kantiner. Det er blevet tradition, at det ved denne hyggelige begivenhed er en yndig, bredskuldret, mandlig FU(AN) der agerer luciabrud.

\subsection{Nye studerende}
Hvert år i august gentager de samme tumultagtige scener sig; ca. 350 forvirrede unge mennesker myldrer for første gang ind på vores dejlige fakultet. For at disse unge mennesker skal føle sig velkomne, afholder \TKET adskillige arrangementer. Det ene er naturligvis SØØL. Første gang kammeret inviterede russerne over til søen var i 1984, og siden er det blevet afholdt hvert eneste år. Hvis regnen står ned, flytter man lidt af søen ind i Vandrehallen og kalder det vandreøl. Men der er nu bedst ved søen. Senere arrangeres også Rushyg, der, som bekendt, er den fest, hvor kammeret viser sig fra sin gavmilde side og giver russerne én øl. Men sådan har det ikke altid været. Allerede i 1957 eksisterede der noget der hed rusmodtagelse, hvor de nye blev præsenteret for stedet og for \TKET, der ved dette arrangement gav øl og smørrebrød. Denne modtagelse havde dog i høj grad et fagligt præg, noget der siden er overtaget af tutorerne, der så lader TK være fest og (én) pils. Siden 1987 har det ligeledes været tradition, at et bestyrelsesmedlem (oftest \CERM) præsenteres til den officielle velkomst som en vigtig person, der har noget vigtigt at fortælle russerne. Herefter forsøger han at binde de stakkels små en grum løgnehistorie på ærmet ved hjælp af bl.a. plantede russer og tutorer. I 1987 startede løjerne med at FORM lagde kortene på bordet, og fortalte russerne, at om et halvt år var deres antal halveret. Det mundede altsammen ud i, at en plantet rus rejste sig og gav FORM en lagkage i hovedet. Den lyvende person afbrydes gerne efter passende tid af resten af BEST/FU, der præsenterer sig og synger en sang om kammeret. De seneste år har det været en sang skrevet på menuetten fra Elverhøj, der første gang blev fremført til et rusteater i 1995 (for tekst: Se jubilæumssangbogen).

\subsection{BODYCRASHING}
Den mærkværdige sport bodycrashing tog sin spæde begyndelse i det tidlige forår 1985, nærmere betegnet i slutningen af februar. Igennem en uges tid havde NæstFormationen og FORMjuntaen gejlet hinanden op ved at udbrede løbesedler skrevet på de respektive foreningers skrivemaskiner, \RemToR og Remington. Det store uundgåelige brag kom så endeligt: Den Dræbende Fløde var skabt. Bodycrashing, som det snart udviklede sig til, er siden blevet kendt ikke bare på universitetet eller i byen. Efterhånden er rygtet om denne bizarre konkurrence nået ud til det ganske land, og således var en journalist fra Jyllands-Posten i 1995 modig nok til at stille op i $2^4 - 3$ og skrive en artikel derom.

\subsection{Brian}
Den ædle personlighed vagtbrien Brian startede sin karriere som et stykke brie, der var blevet efterladt i \TKETs varetægt. Derfor blev han puttet i en plasticpose og hængt op på opslagstavlen. Da han var blevet for flydende, blev han under navnet Brie-ann den 23. september 1987 puttet i en brun glasflaske. Den vogtede så i over to år på \TKET og blev ved fuldmåne fodret med brieost. Men den 28. oktober 1989 blev den på en løvfaldstur begravet, idet den havde fået magtfulde modstandere. Dette er dog for intet at regne når man er en standhaftig brie, og nøjagtig et år senere på en mørk og stormfuld nat (den er god nok), genopstod brien fra selvsamme sted, den var nedgravet. Denne gang lydende navnet Brian. Således gik det til at \TKET igen havde sin vagtbrie på hæderspladsen i vindueskarmen. I 1992 kom der endda en kone til: Man fandt nemlig under renovationen af Fysisk Kantine i en af automaterne et (mindst) ti år gammelt stykke sandkage, der blev stillet ved Brians side. Naturligvis hed dette stykke sandkage Sanne. Én gang i historien, nemlig i efteråret 2002, har Brian fået lejlighed til at vise sin styrke. Der skete nemlig \textit{ikke}(!) det, at en uheldig kammerhænger kom til at puffe ham ned fra den faste plads. Det er så vidt vides hverken før eller siden sket at hele B-bygningen, A-bygningen og store dele af fysik er blevet evakueret! Brian stod ikke til at redde og blev begravet ved Kalø Slotsruin sammen med Sanne, der fulgte ham i døden. Siden er han genopstået som en ny og endnu stærkere Brian.

\subsection{Verdens Kedligste Foredrag}
\TKET er jo blandt andet en foredragsforening, og i 1989 besluttede man, at der burde afholdes en årligt tilbagevendende konkurrence, for at finde ud af hvem der kunne holde „Verdens Kedeligste Foredrag“. Man indbød derfor en „mad dog“ matematiker, en „mad dog“ datalog og en „mad dog“ fysiker, der så hver skulle holde et foredrag af ca. 25 minutters varighed. Vinderen blev begavet med en vandrepokal og retten til at bære titlen „ked.scient.“ indtil der var fundet en ny vinder. Hvis man vinder tre år i træk, får man pokalen og titlen til ejendom. En af de nyeste foredragstraditioner i kammerregi er det årlige stand-up-arrangement. Der indbydes tre stand-up-komikere og aftenen slutter, som øvrige foredrag, med en café.

\subsection{FU'er}
Siden \TKETs stiftelse har vi haft et festudvalg; de hjælpsomme, om end til tider dovne, FU’er. Disse har mangfoldige arbejdsopgaver som fx at fylde køleskabet, slæbe ting og dele løbesedler ud. Den 30. april 1993 afholdtes det første FU-løb, hvor FU’erne dystede på at uddele løbesedler hurtigst, mest talentfuldt, under mest fedten for dommerne og ikke mindst yndigst (for passende definition). Dette løb er nu en årligt tilbagevendende kappestrid om ære og andet mundgodt.

\subsection{Kammerets regler}
Nogle folk mener at \TKET er et vildt forvirrende sted med en masse underlige regler, der bare opretholdes for at få andre til at betale vores øl. Det er kun delvist rigtigt, idet mange af reglerne jo er til for at beskytte folk og inventar. Hvordan får man fx en skrivemaskine til at virke i 50 år hvis man må åbne øl med den? At øloplukkere heller ikke må bruges på \TKET er også klart; de ville bare blive slidt op inden de nåede at blive brugt til en eneste fest. Og desuden ville de bare blive væk. Det kan fortælles, at man i januar 1974 overvejede om stolevipning skulle forbydes, da stole og gulv så ville slides mindre. Det vides ikke om det blev vedtaget, men det er i hvert fald ikke en regel der håndhæves i dag; stolene er dem portnerne allerede har opgivet, og gulvet kaldes jo ikke kammerfælden uden grund. I majrevyen 1987 opførtes en sketch om „Den uafhængige kontrasummeriske fronts ekstremistiske militante fraktion“, der kom på scenen med $\Sigma$-forbudt-skilte. Et af disse skilte er siden blevet hængt op på \TKET og det er således ikke længere tilladt at skrive sumtegn her. Indenfor de seneste år er der indført en række nye regler som følge af kammerpolitiske strømninger. Det er således ikke for tiden tilladt at være nøgen på \TKET, ligesom man ikke må ryge eller have en forstyrrende telefon under bestyrelsesmøder.

\subsection{Inventar}
\TKET har en frygtelig masse inventar, som det naturligvis er ulovligt at hærge. Et af de vigtigste aggregater er køleskabet, så vi kan få kolde øl. Det første køleskab, Majoren, fik vi af et kammerpar, der blev gift, og pludseligt havde et i overskud. Siden fik vi et køleskab af kantinebestyreren, så det kom naturligvis til at hedde Køle-Jørgen, og i dag kører vi på tredje skab, Dunken. I øvrigt kan nævnes at man i 1991–1992 havde et årgangsinventar, nemlig Ole Mørkholt Andersen, der hang så meget på kammeret, at man anså det for rimeligt at erklære ham for inventar. Han måtte således heller ikke hærges det år. Da det blev tradition at \TKET skulle stille op i den panuniversitære kapsejlads på uni-søen, skabtes TK Engineering – en task force af handymen (m/k). Disse innovative mennesker står, udover en række prægtige stridsbåde, også bag kammerets slagbænk, Jar Jar Bænk, som udover at give ekstra siddepladser kan rumme alskens herligheder (fx en FORM).

\subsection{Tavlen}
\TKET har en tavle. Denne tavle har bl.a. et omgangsskyldnerfelt, som er blevet brugt til at notere syndere på, der så har måttet gøre bod. Oprindeligt kunne alle sætte alle „på tavlen“, men da dette til tider førte til tvivlsomme tilfælde, indførte man i 1999 det såkaldte „diskofelt“. Siden da er det kun FORM der må sætte folk på tavlen. Alle andre kan skrive syndere i diskofeltet, og klandringen tages op til diskussion på et senere bestyrelsesmøde. Året efter blev diskofeltet til en diskokugle.

\begin{center}
    \scshape{I øvrigt mener vi at \TKET vil bestå!}
\end{center}
\begin{flushright}
    Jonas Cilieborg, SEKR 1991–1992, \CERM 1992–1993
    
    Redigeret og opdateret af jubilæumsskriftudvalget 2006
    
    Afsnitstitler tilføjet af jubilæumsskriftudvalget 2026
\end{flushright}
